\documentclass[../main.tex]{subfiles}

\begin{document}

\ifSubfilesClassLoaded{

    %\tableofcontents
    
    %\setcounter{chapter}{}
}{}

\chapter{ Quotient Manifolds }
% 代靖涵 25120222201319
\section{Quotients of Manifolds by Group Actions}


\section{Homogeneous Spaces}

\begin{definition}{}{}
     A smooth manifold endowed with a transitive smooth action by a Lie group $G$ is called a \dfntxt{homogeneous $G$-space} .
\end{definition}
\begin{theorem}{Homogeneous Space Construction Theorem}{}
Let $G$ be a Lie group and let $H$ be a closed subgroup of $G$. The left coset space $G/H$ is a topological manifold of dimension equal to $\operatorname{dim} G - \operatorname{dim} H$, and has a unique smooth structure such that the quotient map $\pi : G \to G/H$ is a smooth submersion. The left action of $G$ on $G/H$ given by
\begin{equation}
g_1 \cdot (g_2 H) = (g_1 g_2) H \tag{21.3}
\end{equation}
turns $G/H$ into a homogeneous $G$-space.
\end{theorem}

\begin{theorem}{Homogeneous Space Characterization Theorem}{}
Let $G$ be a Lie group, let $M$ be a homogeneous $G$-space, and let $p$ be any point of $M$. The isotropy group $G_p$ is a closed subgroup of $G$, and the map $F: G/G_p \to M$ defined by $F(gG_p) = g \cdot p$ is an equivariant diffeomorphism.
\end{theorem}

\section{Quotient Lie Groups}
\end{document}